\begin{frame}{Comparação entre teoria e experimento}
  \begin{columns}[T,onlytextwidth]
    \begin{column}{0.48\textwidth}
      \centering
      \textbf{Circuito 1}

      \begin{tikzpicture}
        \begin{semilogxaxis}[
          grafico,xmin=0.3,xmax=50,ymin=0,ymax=10.5,
          xlabel={Frequência (Hz)},ylabel={Magnitude $|H|$ (V/V)},
          legend pos=north east
        ]
          \addplot[very thick,RoxoPrincipal] table[x=f,y expr={10^(\thisrow{ganho1_db}/20)}]{slides/dados/resposta_frequencia.dat};
          \addlegendentry{Modelo nominal}
          \addplot[only marks,mark=*,mark size=2pt,MalvaDestaque]
            table[x=f1,y expr={\thisrow{vo1_v}/\thisrow{vi1_v}}]{slides/dados/medicoes_frequencia.dat};
          \addlegendentry{Bancada}
        \end{semilogxaxis}
      \end{tikzpicture}
    \end{column}
    \begin{column}{0.48\textwidth}
      \centering
      \textbf{Circuito 2}

      \begin{tikzpicture}
        \begin{semilogxaxis}[
          grafico,xmin=1,xmax=200,ymin=0,ymax=1.15,
          xlabel={Frequência (Hz)},ylabel={Magnitude $|H|$ (V/V)},
          legend pos=south west
        ]
          \addplot[very thick,RoxoPrincipal] table[x=f,y expr={10^(\thisrow{ganho2_db}/20)}]{slides/dados/resposta_frequencia.dat};
          \addlegendentry{Modelo nominal}
          \addplot[only marks,mark=*,mark size=2pt,MalvaDestaque]
            table[x=f2,y expr={\thisrow{vo2_v}/\thisrow{vi2_v}}]{slides/dados/medicoes_frequencia.dat};
          \addlegendentry{Bancada}
        \end{semilogxaxis}
      \end{tikzpicture}
    \end{column}
  \end{columns}

  \small
  \begin{block}{Comportamentos observados}
    A atenuação do circuito 1 é gradual. No circuito 2, o ganho permanece próximo
    da unidade e a frequência de pico registrada é $17{,}5\,\mathrm{Hz}$.
  \end{block}
  \centering\scriptsize
  Curvas com componentes nominais; pontos calculados por $|H|=V_o/V_i$
  com as tensões de entrada e saída medidas em cada frequência.
\end{frame}
